%!TEX root = ./../../../thesis-trier.tex

Our first result is a consequence of Theorem~\ref{prelim:thm:sync_words_is_ideal}. Although it has a simple proof,
it is quite important for many of the following results.

\begin{theorem}
	\label{constr_sync:thm:L_in_FactLbar}
	Let $L, L' \subseteq \Sigma^*$ with $L \subseteq \factor{L'}$.
	If $\semiautomat{A}{Q}{\Sigma}{\delta}$ 
	has a synchronizing word in $L$, then it has one in $L'$.
\end{theorem}
\begin{proof}
	Suppose $u \in L$ synchronizes $\mathcal A$.
	By assumption there exist $v, w \in \Sigma^*$
	such that $vuw \in L'$.
	By Theorem~\ref{prelim:thm:sync_words_is_ideal} the
	latter word synchronizes $\mathcal A$ as well.
\end{proof}

Note that $L \subseteq \factor{L'}$ is equivalent to $\factor{L} \subseteq \factor{L'}$.
%As $L \subseteq L'$ implies $\factor{L} \subseteq \factor{L'}$
%and $\factor{\factor{L}} = \factor{L}$, the next is implied.
This gives the following consequence of Theorem~\ref{constr_sync:thm:L_in_FactLbar}.

\begin{corollary}
	\label{constr_sync:cor:factors_sets_included}
	Let $L, L' \subseteq \Sigma^*$ with $\factor{L} = \factor{L'}$.
	Then every automaton $\mathcal A$ has a synchronizing word in $L$
	if and only if it has one in~$L'$.
\end{corollary}

For any $L \subseteq \Sigma^*$ we have
$\factor{L} = \factor{\pref{L}} = \factor{\suff{L}}$. So with Corollary~\ref{constr_sync:cor:factors_sets_included},
we conclude that a given automaton has a synchronizing word in one of
the languages $L$, $\factor{L}$, $\pref{L}$, or $\suff{L}$ if and only if it has one in all of these languages.
As a consequence, we obtain that the corresponding constrained synchronization problems
are equivalent for $\leACZEROm$-reductions (in fact,
the input can be passed through without any modifications
and so these problems are also equivalent
for reductions by projections as considered in~\cite{valiant1982rfducibility}).
%, but $\leACZEROm$-reductions is sufficient for our applications). 

\begin{corollary}
	\label{gen_results:cor:factors_etc_complexity}
	Let $L \subseteq \Sigma^*$ be a regular language. 
	The following problems are equivalent for~$\leACZEROm$-reductions.
	\begin{enumerate}
		\item $\constrsync{L}$
		\item $\constrsync{\factor{L}}$
		\item $\constrsync{\pref{L}}$
		\item $\constrsync{\suff{L}}$
	\end{enumerate}
	Furthermore, if for two languages $L, L'\subseteq \Sigma^*$ the factor, prefix, or suffix sets %closures
	are equal, then $\constrsync{L} \equiv_m^{\AC^0} \constrsync{L'}$.
\end{corollary}

\begin{corollary} \label{constrsync:genresults:cor:union-fact}
	Let $L, L' \subseteq \Sigma^*$ be regular languages.
	Then we have  $\constrsync{L \cup L'} \eqACZEROm \constrsync{\fact{L}\cup L'}$
\end{corollary}
\begin{proof}
	This follows from \cref{gen_results:cor:factors_etc_complexity}
	since $\fact{L \cup L'} = \fact{\fact{L}\cup L'}$.
\end{proof}

\begin{corollary} \label{constrsync:genresults:cor:left-right-ideal}
	Let $L \subseteq \Sigma^*$ be a regular language
	containing a non-empty left or right ideal language.
	Then we have $\constrsync{L} \eqACZEROm \constrsync{\Sigma^*}$.	
\end{corollary}
\begin{proof}
	Without loss of generality, 
	we have $I \subseteq L$
	for some non-empty left ideal language~$I \subseteq \Sigma^*$ (which might be non-regular).
	Pick some $u \in I$ and consider a word $v \in \Sigma^*$.
	Then $vu \in I$
	so that $v \in \fact{I}$.
	This gives $\Sigma^* \subseteq \fact{I} \subseteq \fact{L}$.
	Thus $\fact{L} = \Sigma^*$
	and the claim follows from \cref{gen_results:cor:factors_etc_complexity}.
\end{proof}

\begin{corollary}
	Let $I, J \subseteq \Sigma^*$ be a left and a right ideal, respectively,
	which are both non-empty.
	Then $\constrsync{I\cap J} \eqACZEROm \constrsync{\Sigma^*}$.
\end{corollary}
\begin{proof}
	By assumption, we have $u \in I$ and $v \in J$ for some $u,v \in \Sigma^*$.
	Let $w \in \Sigma^*$.
	Then we have $vwu \in I$ since~$I$ is a left ideal 
	and $vwu \in J$ since~$J$ is a right ideal.
	Thus $vwu \in I \cap J$.
	This gives $\Sigma^* \subseteq \fact{I \cap J}$
	and the claim follows from \cref{gen_results:cor:factors_etc_complexity}. 
\end{proof}